\section{Introduction}

Teleoperation leverages human judgment to perform remote manipulation tasks that remain difficult for autonomous robots, particularly in hazardous environments, flexible manufacturing, and unstructured settings. Haptic feedback conveys remote contact information to the operator and supports assessment of contact state, grasp stability, and slip. Previous studies have examined the stability and transparency of bilateral teleoperation and the adaptation of controllers to the environment, operator, and task \cite{ref01,ref02}. In practical mechatronic systems, however, bilateral control must operate together with perception, mechanical execution, and the human--machine interface.

For heterogeneous objects, a fixed parameter set must trade off gentle contact against positioning and transfer stability. Lower stiffness and grasp force can reduce contact severity but may impair positioning and holding stability, whereas higher settings can improve positional robustness at the cost of greater impact or deformation risk. These competing requirements motivate object-conditioned configuration across multiple subsystems rather than a single fixed setting.

Impedance control regulates the relationship among displacement, velocity, and interaction force, allowing a robot to remain compliant during contact \cite{ref03}. Fixed impedance is straightforward to implement but inevitably represents a compromise across objects. Variable-impedance methods instead adapt stiffness and damping according to task state, contact information, or learned policies, while accounting for the stability implications of time-varying parameters \cite{ref04,ref05}. Passivity-preserving variable-stiffness control has also been addressed with energy-tank formulations that explicitly manage energy exchanged during stiffness variation \cite{ref28}. Previous teleoperation and tele-impedance studies have extensively regulated remote-robot stiffness and impedance according to operator motion, estimated human impedance, contact state, or task phase \cite{ref06,ref07,ref08,ref09,ref18,ref19,ref20}. Bilateral force feedback and operator-facing haptic feedback have also been integrated with impedance-regulated teleoperation \cite{ref08,ref17,ref22,ref25}. These studies show that remote mechanical compliance and operator-facing feedback can be integrated within teleoperation systems, while their primary emphasis is generally placed on impedance transmission, stiffness regulation, transparency, or feedback presentation.

Visual and multimodal information has subsequently been used to support semi-autonomous impedance adaptation. A vision-and-voice tele-impedance method selected slave-side impedance from object properties while retaining an operator confirmation or correction step \cite{ref10}, and visual-impedance control has combined visual and force information in a visual-feature space \cite{ref11}. At the workflow level, intention inference and task-oriented shared control have reduced manual decision-making \cite{ref12}. More recent approaches have derived stiffness from object geometry, material, and environmental relationships \cite{ref13} or generated task-relevant stiffness matrices from visual, gaze, and language inputs \cite{ref14}. Thus, visual object attributes, geometry and material, task intent, and related multimodal inputs have primarily been translated into remote stiffness, impedance matrices, or control-authority decisions.

Related studies on automatic switching and passive arbitration have addressed how control authority is allocated and when it changes \cite{ref15,ref16}, and haptic interfaces have been evaluated for telemanipulation performance \cite{ref17}. Teleoperated grasping work has also used tactile deformation and grip-safety information to regulate gripper impedance for maintaining grip force \cite{ref30}. These lines of work demonstrate the value of object- and task-dependent configuration, but they typically address the remote arm, the operator interface, or the gripper as separate design elements. Among the closely related studies reviewed here, we did not identify a framework in which one detected object-semantic condition directly assigns a common parameter bundle spanning slave-arm impedance, master-side haptic rendering, and gripper execution. The unresolved system-level question is therefore not whether vision can regulate impedance or whether haptic and gripper parameters can be adapted individually, but whether a single interpretable operational strategy can configure these three subsystems jointly and yield task-level benefits beyond visual-semantic display, manual strategy selection, or impedance-only reconfiguration.

To address this question, we propose an object-conditioned strategy-layer configuration architecture designed for contact-preparatory parameter setting in haptic teleoperation of heterogeneous objects. Each detected object class is mapped to an operational strategy and then to a joint parameter vector. After task initiation, the first threshold-qualified mapped detection initiates a single configuration transaction, and the selected strategy is locked for the remainder of the trial. The asynchronous implementation keeps low-rate perception outside the nominal supervisory call path and avoids continuous online optimization. Because operator motion was not gated by perception and synchronized trigger--contact timestamps were not recorded, trial-level contact-preparatory timing could not be verified.

The central positioning of this work is an interpretable, strategy-layer configuration architecture rather than post-contact continuous optimization or persistent visual servoing. Object class serves as a task prior for a coordinated, atomic assignment of the slave-side mechanical response, the operator-facing haptic interface, and gripper execution. This assignment creates a consistent operating state across three otherwise separate mechatronic subsystems; it is not a dynamically derived coupling law, impedance-matching relationship, or statistically demonstrated synergy among channels. The architectural contribution lies in combining semantic strategy selection, bounded asynchronous queues, non-blocking supervisory polling, one-shot event triggering, and intra-trial strategy locking into a constrained configuration workflow. It is therefore a system-level configuration architecture and experimental evaluation, rather than a new impedance controller or visual detector. The formal study evaluates this architecture in a closed set of six known object instances; it does not establish generalization to unseen objects or verify trial-level trigger correctness.

The system-level evaluation uses five modes: fixed-parameter baseline (A), manual-selection workflow baseline (B), joint parameter-bundle configuration (C), visual-semantic display with fixed parameters (D), and impedance-only reconfiguration (E). The evaluation focuses primarily on whether the complete parameter bundle improves task completion time relative to impedance-only reconfiguration under matched visual-semantic and impedance conditions. The C--D comparison assesses joint parameter-bundle configuration beyond the visual-semantic display, while trajectory length, pause count, success rate, and workload are interpreted as secondary or exploratory outcomes. Operator- and object-stratified results describe whether this difference has a consistent direction; they do not identify the independent causal contributions of the haptic and gripper parameter groups.

The main contributions are as follows:

\begin{enumerate}
\item \textbf{Object-conditioned strategy-layer parameter-bundle architecture.} The framework elevates object class from a display cue to a strategy-layer event that atomically assigns a seven-dimensional parameter bundle across slave mechanical response, operator-facing haptic rendering, and gripper execution. Object classes are grouped into fragility-priority, balanced, and stability-priority operational strategies and mapped to operating points constrained by hardware capabilities, task risks, and preliminary screening. The contribution is consistent configuration across subsystems, not a dynamic coupling model between channels.
\item \textbf{Asynchronous, event-triggered one-shot reconfiguration.} Image acquisition at 15 fps and independent YOLO11n inference communicate with a nominal 200-Hz supervisory update through bounded queues. Non-blocking polling isolates perception from the supervisory call path; one-shot triggering converts semantic output into a single configuration transaction; strategy locking avoids repeated intra-trial switching caused by detection fluctuations; and an approximately 300-ms transition smooths changes in stiffness and the controller damping-scaling parameter.
\item \textbf{Human-in-the-loop system-level evaluation across five modes.} Five modes were implemented on an Omega.7--Panda platform to distinguish visual-semantic display with fixed parameters, a manual-selection workflow baseline, impedance-only reconfiguration, and joint parameter-bundle configuration. The primary C--E comparison evaluates the complete parameter bundle at the system level and does not separate the individual effects of the added haptic-interface and gripper channels. The evaluation involved three operators, six objects, and 135 trials; its inferential scope is limited to the repeated measurements in this sample.
\end{enumerate}
